\hypertarget{v11-coupled_lab_full-exact-laboratory-and-certified-evaluation}{%
\chapter[Exact laboratory and certified evaluation]{Exact laboratory and certified
evaluation}\label{app:laboratory}\label{v11-coupled_lab_full-exact-laboratory-and-certified-evaluation}}

The reusable module \path{coupled_lab.py} uses Python's standard
library and imports the existing exact coefficient, phase and tap
operations from \path{family_lab.py}. The figure builder additionally
uses NumPy and Matplotlib. Neither script contacts an external service.
The generic integer and signed catalog rows are formal power series,
with no claim of analytic convergence.

\hypertarget{v11-coupled_lab_full-normalization-and-recurrence}{%
\section{Normalization and
recurrence}\label{v11-coupled_lab_full-normalization-and-recurrence}}

Write

\[
U(z)=\sum_{n\ge0}A_n\frac{z^{2n+1}}{(2n+1)!},\qquad
V(z)=\sum_{n\ge0}B_n\frac{z^{2n+1}}{(2n+1)!}.
\]

The exact recurrence is

\[
A_n=\sum_{i+j=n-1}\binom{2n}{2i+1}q^{2i+1}r^{2j+1}A_iB_j,
\qquad
B_n=\sum_{i+j=n-1}\binom{2n}{2i+1}a^{2i+1}b^{2j+1}A_iB_j,
\]

with \(A_0=B_0=1\). Coefficients with even derivative index vanish.
Swapping the components sends the parameter tuple \((q,r,a,b)\) to
\((b,a,r,q)\). Identical rows \((a,b)=(q,r)\), or exchanged rows
\((a,b)=(r,q)\), yield the scalar reduction \(U=V=T_{q,r}\).

Both derivative arrays are polynomials in four scale variables with
nonnegative integer coefficients. Their total degree is at most
\(n(n+1)\). Indeed, a recurrence term with \(i+j=n-1\) has degree at
most \(2n+i(i+1)+j(j+1)\le n(n+1)\). If all four denominators divide
\(D\), then \(D^{n(n+1)}A_n\) and \(D^{n(n+1)}B_n\) are integers. The
clearing factor is a convenient uniform choice and need not be minimal.

\hypertarget{v11-coupled_lab_full-an-elementary-exact-evaluation-certificate}{%
\section{An elementary exact evaluation
certificate}\label{v11-coupled_lab_full-an-elementary-exact-evaluation-certificate}}

This argument uses \(0<q,r,a,b<1\) and does not require the order
theorem. Let \(Q=\max(q,r,a,b)\). All ordinary odd Taylor coefficients
are positive and at most one: induction in the ordinary recurrence
bounds the coefficient of \(z^{2n+1}\) by \(n/(2n+1)\) for \(n\ge1\).
Thus both series converge on the unit disk and, for \(0\le x<1\),

\[
\max(U(x),V(x))\le \frac{x}{1-x^2}.
\]

Analytic continuation by the integral equations expands each known disk
by the factor \(1/Q\), so both solutions are entire. Positive
coefficients give \(\max_{|z|\le x}|F(z)|=F(x)\) for either component.
The integral equations therefore imply

\[
M(x):=\max(U(x),V(x))\le x\bigl(1+M(Qx)^2\bigr).
\]

Start at the first member \(Q^kx<1\) of this finite contraction ladder,
use \(x/(1-x^2)\) there, and repeatedly apply \(B(x)=x(1+B(Qx)^2)\)
backward. This defines the rational \path{disk_majorant} bound for
rational parameters and radius. The implementation has explicit
recursion and integer size guards, which raise errors rather than
returning a weakened result.

If \(F_N\) retains odd indices \(0,\ldots,N\), then for \(S>R\ge0\),

\[
\max_{|z|\le R}|F(z)-F_N(z)|
\le B(S)\left(\frac RS\right)^{2N+3}.
\]

For every omitted exponent \(m\ge2N+3\), use \(R^m\le (R/S)^{2N+3}S^m\),
sum positive terms, and bound their sum by \(F(S)\le B(S)\). No extra
geometric denominator is needed. For rational \(x\ge0\),
\path{certified_evaluation} returns the exact interval
\([F_N(x),F_N(x)+B(S)(x/S)^{2N+3}]\) for each component.

\Needspace{23\baselineskip}
\hypertarget{v11-coupled_lab_full-worked-calls}{%
\section{Worked calls}\label{v11-coupled_lab_full-worked-calls}}

\begin{verbatim}
from fractions import Fraction as Q
import coupled_lab as lab

p = (Q(1,5), Q(2,5), Q(3,5), Q(7,10))
A, B = lab.coupled_odd_derivatives(*p, last=12)
certificate = lab.certified_evaluation(*p, Q(1), last=30, outer_radius=2)
root = lab.pf_order(*p)  # floating approximation, not interval arithmetic

derivatives = lab.full_egf_coefficients(B)
phases = lab.phase_lanes(derivatives, 3)
matrix = lab.polyphase_tap_matrix([1, -2, 1], 3)
length = len(derivatives) + 2
output_phases = lab.apply_polyphase_matrix(phases, matrix, length)
output = lab.merge_lanes(output_phases, length)
\end{verbatim}

At \(x=1\), with \(N=30\) and \(S=2\), the common tail bound is below
\(1.600\times10^{-13}\). The partial values are approximately
\(U(1)=1.027044341827337\) and \(V(1)=1.1467074957737624\). These
displayed decimals are rounded illustrations. The returned fractions are
the actual certificate.

Run the independent checks and regenerate figures from the package\textquotesingle s \texttt{research/} directory:

\begin{verbatim}
python3 checks/coupled_v11.py
python3 build_coupled_figures.py
\end{verbatim}

The first command passes 1,559 exact assertions, 10 numerical
comparisons, and 13 separate input/resource guards. Independent
constructions include dense ordinary Taylor Picard iteration,
sine/cosine division for tangent, the exact sinh formula, component
symmetry, and direct convolution versus phase carry. The test counts
describe bounded implementation checks, not universal proofs.
